\documentclass[12pt,a4paper]{article}

% ==================================================
% PACKAGES
% ==================================================
\usepackage[utf8]{inputenc}
\usepackage[T5]{fontenc}
\usepackage[vietnamese]{babel}
\usepackage{geometry}
\usepackage{array}
\usepackage{longtable}[=2020-01-02]
\makeatletter
\long\def\LT@p@ftntext#1{%
    \edef\@tempa{%
        \the\LT@p@ftn
        \noexpand\footnotetext[\the\c@footnote]%
    }%
    \global\LT@p@ftn\expandafter{\@tempa{#1}}%
}
\makeatother
\usepackage{booktabs}
\usepackage{tabularx}
\usepackage{ragged2e}
\usepackage{float}
\usepackage{enumitem}
\usepackage{xcolor}
\usepackage{colortbl}
\usepackage{fancyhdr}
\usepackage[most]{tcolorbox}
\usepackage{microtype}
\usepackage{hyperref}

\geometry{left=2cm,right=2cm,top=1.8cm,bottom=1.8cm}
\setlength{\parindent}{0pt}
\setlength{\parskip}{0.35em}

% ==================================================
% STYLE
% ==================================================
\definecolor{VinBlue}{RGB}{0,56,101}
\definecolor{VinGold}{RGB}{190,145,50}
\definecolor{MondBlue}{HTML}{173B57}
\definecolor{MondTeal}{HTML}{087E8B}
\definecolor{MondLight}{HTML}{F2F7F8}
\definecolor{MondRed}{HTML}{B54747}
\definecolor{MondGray}{HTML}{5B6770}

\newcolumntype{Y}{>{\RaggedRight\arraybackslash}X}
\newcolumntype{P}[1]{>{\RaggedRight\arraybackslash}p{#1}}
\newcolumntype{C}[1]{>{\Centering\arraybackslash}p{#1}}

\newtcolorbox{summarybox}[1][]{
    colback=MondLight,
    colframe=MondTeal,
    boxrule=0.8pt,
    arc=2mm,
    title=#1,
    fonttitle=\bfseries
}

\hypersetup{colorlinks=true,linkcolor=VinBlue,urlcolor=VinBlue}

\pagestyle{fancy}
\fancyhf{}
\fancyhead[L]{\small\color{VinBlue}\textbf{MondAIy -- Team Nerd}}
\fancyhead[R]{\small\color{VinBlue}\textbf{Nhật ký làm việc}}
\fancyfoot[C]{\small\color{VinBlue}\thepage}
\renewcommand{\headrulewidth}{0.6pt}
\setlength{\headheight}{15pt}

\begin{document}
\setlength{\emergencystretch}{3em}
\hbadness=10000
\vbadness=10000
\hfuzz=20pt
\vfuzz=20pt
\AddToHook{para/begin}{\hbadness=10000\hfuzz=\maxdimen}

% ==================================================
% TITLE AND TEAM STRUCTURE
% ==================================================
\begin{center}
    {\LARGE\color{VinBlue}\textbf{NHẬT KÝ LÀM VIỆC THEO TUẦN}}\\[0.35em]
    {\Large\textbf{Dự án MondAIy}}\\[0.3em]
    {\color{VinGold}\rule{8cm}{1pt}}
\end{center}

\section{CƠ CẤU NHÂN SỰ VÀ PHÂN CÔNG}

\begin{summarybox}[Nguyên tắc phân công]
Cả bốn thành viên đều đảm nhiệm vai trò \textbf{AI Engineer}. Mỗi người đồng thời
sở hữu một vai trò bổ sung để bảo đảm dự án có người chịu trách nhiệm rõ ràng về
tiến độ, sản phẩm, trải nghiệm, chất lượng và kiến trúc kỹ thuật. Công việc được
đánh giá theo artifact và kết quả bàn giao, không dựa đơn thuần trên số commit hoặc
số dòng mã.
\end{summarybox}

\begin{table}[H]
\centering
\caption{Cơ cấu nhân sự và phạm vi trách nhiệm}
\small
\sloppy
\renewcommand{\arraystretch}{1.22}
\begin{tabularx}{\textwidth}{P{2.4cm}P{3.6cm}Y P{4.3cm}}
\rowcolor{MondBlue}
\color{white}\textbf{Thành viên} & \color{white}\textbf{Vai trò bổ sung} &
\color{white}\textbf{Trách nhiệm chính} & \color{white}\textbf{Artifact tiêu biểu}\\
\toprule
Thái Hoài An & Project Manager (PM) & Điều phối kế hoạch, theo dõi Jira, tổ chức
workflow và hỗ trợ triển khai các module Agent & Jira plan, Discord ingestion,
Telegram, Meeting Minutes, Memory, Google Drive search, video trailer\\
\rowcolor{MondLight}
Nguyễn Văn Phong & Technical Lead & Kiến trúc tích hợp, connector, extension,
scanning và độ tin cậy kỹ thuật & Connector Gmail, Calendar, Jira và Slack; browser
extension, screen scanning, Scan-to-Exact-Task và Golden Dataset\\
Lê Thị Yến Nhi & UI/UX Designer & Thiết kế trải nghiệm, cải tiến giao diện theo
feedback và hỗ trợ chất lượng đầu ra Agent & UI/UX draft và mockup, design iteration,
task metadata filtering, Agent response improvement, Pitch Deck\\
\rowcolor{MondLight}
Nguyễn Đức Dũng & Product Owner (PO) và Quality Control (QC) & Định hướng giá trị,
PRD, tiêu chí nghiệm thu, guardrail, dữ liệu kiểm thử và kiểm soát chất lượng & PRD,
concept, test cases, guardrail, intent/task extraction, file analysis, mobile app,
specification và hồ sơ doanh nghiệp\\
\bottomrule
\end{tabularx}
\end{table}

% ==================================================
% WEEKLY LOG
% ==================================================
\clearpage
\section{NHẬT KÝ TRIỂN KHAI 6 TUẦN}

\small
\renewcommand{\arraystretch}{1.2}
\setlength{\tabcolsep}{5pt}
\begin{longtable}{@{}C{1.2cm}P{3.0cm}P{11.8cm}@{}}
\caption{Nhật ký làm việc theo tuần của nhóm dự án MondAIy}
\label{tab:weekly-log}\\
\rowcolor{MondBlue}
\color{white}\textbf{Tuần} & \color{white}\textbf{Thành viên} &
\color{white}\textbf{Công việc thực hiện}\\
\toprule
\endfirsthead

\rowcolor{MondBlue}
\color{white}\textbf{Tuần} & \color{white}\textbf{Thành viên} &
\color{white}\textbf{Công việc thực hiện}\\
\toprule
\endhead

\bottomrule
\endlastfoot

1 & An & Tham gia workshop chốt ý tưởng, phạm vi và kế hoạch phối hợp; tổng hợp đầu việc ban đầu cho nhóm.\\
\rowcolor{MondLight}
1 & Phong & Tham gia chốt ý tưởng; đánh giá sơ bộ tính khả thi kỹ thuật và các hướng tích hợp.\\
1 & Nhi & Tham gia chốt ý tưởng; phác thảo nhu cầu người dùng và định hướng trải nghiệm.\\
\rowcolor{MondLight}
1 & Dũng & Tham gia chốt ý tưởng; làm rõ vấn đề, tuyên bố giá trị và tiêu chí đầu ra ban đầu.\\

2 & An & Thiết lập và quản lý kế hoạch trên Jira; phân rã milestone và đầu việc theo tuần.\\
\rowcolor{MondLight}
2 & Phong & Khảo sát các giải pháp hiện có; xác định khoảng trống kỹ thuật và hướng connector.\\
2 & Nhi & Nghiên cứu người dùng; xây dựng bản draft UI/UX và luồng màn hình chính.\\
\rowcolor{MondLight}
2 & Dũng & Xây dựng PRD, concept và định hướng giá trị; xác định yêu cầu đầu ra cho các module.\\

3 & An & Xây dựng luồng thu thập và chuẩn hóa dữ liệu từ Discord; phối hợp kiểm tra integration.\\
\rowcolor{MondLight}
3 & Phong & Xây dựng connector với Gmail, Google Calendar và Jira; chuẩn hóa interface tích hợp.\\
3 & Nhi & Phát triển UI/UX mockup; hoàn thiện cấu trúc điều hướng và các màn hình cốt lõi.\\
\rowcolor{MondLight}
3 & Dũng & Xây dựng dữ liệu kiểm thử cho luồng lấy dữ liệu và task; thiết kế guardrail cho hệ thống.\\

4 & An & Xây dựng chức năng chat trên Telegram; phát triển tính năng Meeting Minutes.\\
\rowcolor{MondLight}
4 & Phong & Xây dựng browser extension; phát triển connector với Slack và rà soát luồng tích hợp.\\
4 & Nhi & Cải thiện UI/UX dựa trên feedback khách hàng; chuẩn hóa các trạng thái tương tác.\\
\rowcolor{MondLight}
4 & Dũng & Cải tiến Intent Recognition; cải tiến luồng tóm tắt và trích xuất task có bằng chứng.\\

5 & An & Xây dựng cơ chế Memory cho Agent; xây dựng video trailer giới thiệu sản phẩm.\\
\rowcolor{MondLight}
5 & Phong & Xây dựng luồng scanning màn hình; khảo sát khả năng hỗ trợ Zalo và nền tảng bổ sung.\\
5 & Nhi & Cải tiến luồng tóm tắt; xây dựng Metadata Filtering cho task và hỗ trợ kiểm tra output.\\
\rowcolor{MondLight}
5 & Dũng & Xây dựng luồng phân tích file; phát triển mobile app; triển khai nội dung truyền thông fanpage.\\

6 & An & Rà soát và cải tiến hệ thống; xây dựng luồng tìm kiếm trên Google Drive; tổng hợp tiến độ.\\
\rowcolor{MondLight}
6 & Phong & Rà soát hệ thống; cải tiến Scan-to-Exact-Task; xây dựng Golden Dataset và regression cases.\\
6 & Nhi & Rà soát UI; cải thiện phản hồi Agent; hoàn thiện Pitch Deck và tài liệu trình bày sản phẩm.\\
\rowcolor{MondLight}
6 & Dũng & Hoàn thiện hồ sơ doanh nghiệp và specification; thiết kế, kiểm thử và cải tiến theo kết quả testing.\\

\end{longtable}

% ==================================================
% INDIVIDUAL ASSESSMENT
% ==================================================
\section{ĐÁNH GIÁ THEO THÀNH VIÊN}

\setlength{\tabcolsep}{4pt}
\begin{longtable}{@{}P{2.4cm}P{4.5cm}P{4.5cm}P{4.0cm}@{}}
\caption{Điểm mạnh, điểm cần cải thiện và giá trị đóng góp}
\label{tab:member-assessment}\\
\rowcolor{MondBlue}
\color{white}\textbf{Thành viên} & \color{white}\textbf{Điểm mạnh} &
\color{white}\textbf{Điểm cần cải thiện} & \color{white}\textbf{Giá trị nổi bật}\\
\toprule
\endfirsthead
\rowcolor{MondBlue}
\color{white}\textbf{Thành viên} & \color{white}\textbf{Điểm mạnh} &
\color{white}\textbf{Điểm cần cải thiện} & \color{white}\textbf{Giá trị nổi bật}\\
\toprule
\endhead
\bottomrule
\endlastfoot

Thái Hoài An &
Khả năng điều phối và chuyển nhanh giữa quản lý dự án với triển khai AI; bao phủ
Discord, Telegram, Meeting Minutes, Memory và Drive search; chủ động hỗ trợ truyền thông. &
Cần khóa acceptance criteria và dependency rõ hơn trên Jira; tránh để PM trở thành
nút thắt khi đồng thời sở hữu nhiều module; tăng báo cáo định lượng theo milestone. &
Kết nối kế hoạch nhóm với các workflow chủ động và giữ nhịp triển khai xuyên sáu tuần.\\

\rowcolor{MondLight}
Nguyễn Văn Phong &
Năng lực tích hợp hệ thống mạnh; xử lý nhiều connector và nền tảng; có tư duy
technical lead về interface, scanning, test data và độ tin cậy. &
Cần chuẩn hóa ADR, SLO và ownership bảo trì connector; hạn chế mở rộng nền tảng
quá sớm; tăng handover để tri thức kỹ thuật không tập trung vào một người. &
Xây nền kết nối và pipeline kỹ thuật để sản phẩm có thể đọc dữ liệu thực thay vì chỉ demo.\\

Lê Thị Yến Nhi &
Khả năng chuyển feedback thành cải tiến trải nghiệm; kết nối UI/UX với metadata,
chất lượng phản hồi Agent và cách trình bày giá trị sản phẩm. &
Cần gắn mỗi thay đổi thiết kế với usability metric; hoàn thiện design system và
research log; làm rõ ranh giới giữa artifact thiết kế và artifact AI engineering. &
Bảo đảm sản phẩm có trải nghiệm sử dụng nhất quán, dễ trình bày và bám sát nhu cầu người dùng.\\

\rowcolor{MondLight}
Nguyễn Đức Dũng &
Khả năng bao quát sản phẩm và chất lượng; kết nối PRD, dữ liệu kiểm thử, guardrail,
intent/task extraction, tài liệu và kiểm thử end-to-end. &
Cần phân quyền bớt PO/QC/engineering để tránh quá tải; duy trì traceability từ PRD
đến test và release gate; tách rõ prototype, benchmark và production claim. &
Giữ tính nhất quán giữa giá trị sản phẩm, yêu cầu kỹ thuật, bằng chứng kiểm thử và hồ sơ dự án.\\
\end{longtable}

% ==================================================
% QUANTITATIVE CONTRIBUTION
% ==================================================
\section{ĐỊNH LƯỢNG MỨC ĐỘ ĐÓNG GÓP}

Mức đóng góp được chuẩn hóa thành \textbf{100 đơn vị đóng góp}. Việc phân bổ phản
ánh sự bổ trợ giữa các vai trò: người thiên về quản lý/sản phẩm không bị đánh giá
thấp hơn người viết nhiều mã, và artifact thiết kế, kiểm thử, tài liệu hay truyền
thông được ghi nhận khi phục vụ trực tiếp cho kết quả dự án.

\begin{table}[H]
\centering
\caption{Phân bổ đơn vị đóng góp theo nhóm công việc}
\label{tab:contribution-allocation}
\small
\renewcommand{\arraystretch}{1.22}
\setlength{\tabcolsep}{4pt}
\begin{tabular}{P{2.5cm}*{5}{>{\centering\arraybackslash}p{1.75cm}}>{\centering\arraybackslash}p{1.8cm}}
\rowcolor{MondBlue}
\color{white}\textbf{Thành viên} & \color{white}\textbf{\shortstack{Product\\PM}} &
\color{white}\textbf{\shortstack{AI\\Backend}} & \color{white}\textbf{UI/UX} &
\color{white}\textbf{\shortstack{QA\\Data\\Docs}} & \color{white}\textbf{\shortstack{Hỗ trợ\\TT}} &
\color{white}\textbf{Tổng}\\
\toprule
An & 7 & 10 & 1 & 3 & 4 & \textbf{25}\\
\rowcolor{MondLight}
Phong & 3 & 14 & 1 & 6 & 1 & \textbf{25}\\
Nhi & 3 & 6 & 10 & 4 & 2 & \textbf{25}\\
\rowcolor{MondLight}
Dũng & 9 & 8 & 2 & 5 & 1 & \textbf{25}\\
\midrule
\textbf{Toàn nhóm} & \textbf{22} & \textbf{38} & \textbf{14} & \textbf{18} &
\textbf{8} & \textbf{100}\\
\bottomrule
\end{tabular}
\end{table}

\begin{table}[H]
\centering
\caption{Tỷ lệ đóng góp cuối cùng}
\renewcommand{\arraystretch}{1.3}
\begin{tabularx}{0.82\textwidth}{Y>{\centering\arraybackslash}P{3.2cm}>{\centering\arraybackslash}P{3.2cm}}
\rowcolor{MondTeal!15}
\textbf{Thành viên} & \textbf{Đơn vị đóng góp} & \textbf{Tỷ lệ}\\
\toprule
Thái Hoài An & 25/100 & \textbf{25\%}\\
\rowcolor{MondLight}
Nguyễn Văn Phong & 25/100 & \textbf{25\%}\\
Lê Thị Yến Nhi & 25/100 & \textbf{25\%}\\
\rowcolor{MondLight}
Nguyễn Đức Dũng & 25/100 & \textbf{25\%}\\
\midrule
\textbf{Tổng} & \textbf{100/100} & \textbf{100\%}\\
\bottomrule
\end{tabularx}
\end{table}

\section{TỔNG QUAN TIẾN ĐỘ}

Trong sáu tuần triển khai, nhóm phát triển MondAIy theo hướng hoàn thiện đồng thời
sản phẩm, hệ thống AI và trải nghiệm người dùng. Công việc đi từ xác định vấn đề,
PRD và UI/UX đến connector, Agent, Memory, Intent Recognition, Task Extraction,
Guardrail, Scanning và tìm kiếm dữ liệu. Song song, nhóm xây dựng test cases,
Golden Dataset, tài liệu đặc tả, hồ sơ doanh nghiệp, Pitch Deck và nội dung truyền
thông. Kết quả nổi bật không chỉ là số lượng tính năng mà là việc hình thành một
workflow xuyên nguồn có bằng chứng, cơ chế kiểm soát và nền tảng để tiếp tục đánh
giá chất lượng trong các vòng pilot tiếp theo.

\end{document}